\section{Strong irreducibility from local monodromy}\label{sec:irreducibility}
\paragraph{Conventions and inputs.}
Fix a coefficient prime $\ell\ne3$. Write $G_\Q=\Gal(\overline\Q/\Q)$ and
$\tau_\ell$ for the character of $\Q_\ell(-1)$, so that
$\tau_\ell(\Frob_p)=p$ for geometric Frobenius at $p\ne\ell$.
To avoid opposite Hodge--Tate sign conventions, call $h$ a \emph{Tate
exponent} when a Hodge--Tate summand is $\C_\ell(-h)$.
Thus $\Q_\ell(-1)$ has Tate exponent $1$.

The inputs for $V=W_{4,\ell}^+$ are:
\begin{enumerate}
\item $V$ has dimension four over $\Q_\ell$, is unramified outside a finite
set, and is de Rham at $\ell$.
\item Its Tate exponents are $\{0,0,5,5\}$. These follow from the specified
smooth projective Scholl realization of
Theorem~\ref{thm:scholl-realization}, its rational projector and the Hodge
numbers $h^{5,0}=h^{0,5}=2$, using de Rham and Hodge--Tate comparison.
\item At $3$, inertia is unipotent with a monodromy logarithm $N$ having
one Jordan block of length four. With $F$ geometric Frobenius,
\[
 FNF^{-1}=3^{-1}N,
\]
and the eigenvalues on the successive quotients of
$0\subset\ker N\subset\ker N^2\subset\ker N^3\subset\ker N^4=V$
are $3,9,27,81$, respectively.
\end{enumerate}
The third input is the full arithmetic conclusion of
Section~\ref{sec:local-three}, not an inference from the conductor alone. The same construction gives
$W_{4,\ell}^-=W_{4,\ell}^+\otimes\chi_{-3}$.

\begin{lemma}[Rank-one characters over $\Q$]\label{irr:rankone}
Let $\delta:G_\Q\to\Q_\ell^\times$ be a continuous character, unramified
outside a finite set, and Hodge--Tate at $\ell$ with Tate exponent $m$.
Then $\delta=\epsilon\tau_\ell^m$ for a finite-order character $\epsilon$.
\end{lemma}
\begin{proof}
Tate's theorem identifies a Hodge--Tate abelian local representation with
a locally algebraic one when its inertia action is semisimple. A character
automatically satisfies the latter condition. The precise primary reference
is Serre \cite[Chapter~III, \S1.2, Theorem (Tate), printed p.~III-7;
proof in the Appendix to Chapter~III]{SerreRankOne}; the definition in \S1.1
identifies the local algebraic power. In the present $\Q_\ell$-valued
rank-one case, it follows that $\delta\tau_\ell^{-m}$ has finite image
on inertia at $\ell$.

Put $\epsilon=\delta\tau_\ell^{-m}$. By Kronecker--Weber, an abelian
character of $G_\Q$ is a continuous character of
$\widehat\Z^\times=\prod_p\Z_p^\times$. Outside a finite set its
restrictions to the unit factors are trivial, by the unramified hypothesis.
For each remaining $p\ne\ell$, its restriction to $\Z_p^\times$ has
finite image: its open pro-$p$ subgroup maps into $\Q_\ell^\times$,
whose compact subgroups have an open pro-$\ell$ subgroup, so that image
is finite. At $\ell$ finiteness was just proved. Thus the entire image is
contained in the product of finitely many finite groups, and is finite.
The exponent $m$ is the stated Tate exponent, since finite characters
have Tate exponent zero.
\end{proof}

\begin{lemma}[The single-block flag]\label{irr:flag}
Let $k$ be a field of characteristic zero and let $N$ on $k^4$ have one
nilpotent Jordan block. Its invariant subspaces are exactly
$\ker N^r$ for $0\le r\le4$.
\end{lemma}
\begin{proof}
As a $k[t]$-module, with $t$ acting as $N$, this space is
$k[t]/(t^4)$. Its submodules are the ideals $(t^j)/(t^4)$,
equivalently the displayed kernels. This proof remains valid after
arbitrary scalar extension.
\end{proof}

\begin{theorem}\label{irr:absolute}
Under the stated inputs, both $W_{4,\ell}^+$ and $W_{4,\ell}^-$ are
absolutely irreducible for every $\ell\ne3$.
\end{theorem}
\begin{proof}
It suffices to treat $V=W_{4,\ell}^+$. Suppose
$0\ne U\subsetneq V\otimes_{\Q_\ell}\overline\Q_\ell$ is a
$G_\Q$-stable subspace, of dimension $r\in\{1,2,3\}$.
Its finitely many basis coordinates belong to a finite coefficient
extension, so there is no issue of defining a continuous representation
over an infinite coefficient field. There is in fact a stronger descent
argument in this case.

Some inertia element at $3$ acts as $\exp(cN)$ with $c\ne0$.
Since its logarithm is a finite polynomial in $\exp(cN)-1$, stability
under that element makes $U$ stable under $N$. Lemma~\ref{irr:flag}
gives
\[
 U=(\ker N^r)\otimes_{\Q_\ell}\overline\Q_\ell.
\]
Set $U_0=\ker N^r\subset V$. This is defined over $\Q_\ell$ and is
$G_\Q$-stable: extension of scalars makes every $\rho(g)$ preserve it,
and the original matrices are $\Q_\ell$-linear. Hence its determinant
$\delta=\det(U_0)$ is a $\Q_\ell$-valued character.

The category of de Rham representations is closed under subobjects,
quotients and exterior powers; its comparison functor is exact and
strict for the filtrations. Thus $U_0$ and $\delta$ are de Rham, the Tate
exponents of $U_0$ form a submultiset of $\{0,0,5,5\}$, and the
Tate exponent $m$ of $\delta$ is their sum. In particular
\[
 m\in5\Z.
\]
For clarity, the dimension part of closure under subobjects follows already
from left exactness of $D_{\mathrm{dR}}$ and
$\dim D_{\mathrm{dR}}(M)\le\dim M$: apply this to $U_0$, $V$ and
$V/U_0$, with equality for $V$. The usual strictness of de Rham
comparison supplies the asserted filtered submultiset statement.

The character $\delta$ is unramified outside a finite set, so
Lemma~\ref{irr:rankone} gives $\delta=\epsilon\tau_\ell^m$.
Since $U_0$ has unipotent inertia at $3$, its determinant is unramified
there, as is $\epsilon$. On the other hand the explicitly ordered
monodromy flag gives
\[
 \delta(\Frob_3)=3^{1+2+\cdots+r}=3^{r(r+1)/2}.
\]
Therefore
\[
 \epsilon(\Frob_3)=3^{r(r+1)/2-m}.
\]
The left side is a root of unity. A rational integer power of $3$
is a root of unity only when its exponent is zero. Consequently
$m=r(r+1)/2$. For $r=1,2,3$ these values are $1,3,6$, respectively,
none of which belongs to $5\Z$. This contradiction proves absolute
irreducibility of $V$. Tensoring with the one-dimensional character
$\chi_{-3}$ preserves absolute irreducibility and proves the minus case.
\end{proof}

\begin{corollary}[Restriction to finite extensions]\label{irr:strong}
For every finite extension $L/\Q$, the restrictions of either
$W_{4,\ell}^\pm$ to $G_L$ are absolutely irreducible, for $\ell\ne3$.
\end{corollary}
\begin{proof}
It suffices again to treat the plus representation. If a proper nonzero
subspace $U$ is stable under $G_L$, replace $L$ by a finite Galois
extension containing it and put $H=G_L\triangleleft G_\Q$.
The intersection $H\cap I_3$ has finite index in $I_3$, hence its image
under the nonzero tame monodromy parameter still contains a nonzero
element. Consequently every $H$-stable subspace is $N$-stable.
Normality makes every $gU$ also $H$-stable, and all have dimension $r$.
By Lemma~\ref{irr:flag}, all equal $\ker N^r$ after scalar extension.
Thus $U$ is $G_\Q$-stable, contradicting Theorem~\ref{irr:absolute}.
\end{proof}

\paragraph{Limits and consequences.}
The argument uses the written Scholl/Hodge and local-at-$3$ results as
inputs; it does not re-prove their constructions. It makes no assertion
for the coefficient prime $\ell=3$, since the local monodromy calculation
above uses $\ell\ne3$. No conjecture of semisimplicity is used.
Proper induction from a smaller-dimensional representation over a finite
extension is excluded by Corollary~\ref{irr:strong}. An ordinary tensor
product of two rank-two continuous $\ell$-adic representations over finite
coefficient extensions is also excluded:
after a finite local extension its monodromy operator would have form
$N_1\otimes1+1\otimes N_2$, with $N_i^2=0$, hence its cube would be
zero, whereas the actual $N^3$ is nonzero. These exclusions are independent
of the dyadic calculation in Section~\ref{sec:local-two} and do not
establish automorphy or determine the Sato--Tate group.

